\begin{theorem}[The translation component]
\label{thm:c9-P1}
Let $\hbar F_1^{(0)}$ contain the finite counterterm of Theorem \ref{thm:c4-finite} and let $\psi=0$.
\begin{itemize}
\item (i) The quadratic part of the arity-one obstruction along $P_b$, the sum of the coefficients of $\varepsilon^0$ of the defect cycles with two legs and the vertex $S^{(1)}(P_b)$ at the special edge or on a line, vanishes. The quadratic part of the arity-one equation along $P_b$ is therefore solved exactly by
\[
\hbar F_1^{(1)}(P_b)=\text{units}\cdot\Big(\theta_P\,S^{(1)}(P_b)_\perp+\theta_{P,\lambda}\,S^{(1)}(P_b)_\lambda+\theta_{P,V}\,S^{(1)}(P_b)_V\Big)+Q\Big(\kappa'_V\int(\iota_bA|_V,c^*)+\kappa'_\lambda\int(\rho(b)\lambda^*,\lambda^*)\Big),
\]
with the sector parts of Definition \ref{def:c11-kernels} and free numbers $\theta_P,\theta_{P,\lambda},\theta_{P,V},\kappa'_V,\kappa'_\lambda$; in the sector $(A|_{W_{\mathbb R}^\perp},A^*|_{W_{\mathbb R}^\perp})$ it is $\theta_P$ units times $\int(\partial_b(A|_{W_{\mathbb R}^\perp}),A^*|_{W_{\mathbb R}^\perp})$.
\item (ii) For a scalar bilinear $O$, on the part with two legs $A|_{W_{\mathbb R}^\perp}$, the master formula of Theorem \ref{thm:c11-classes} (iii) gives $P_1(O)=\Upsilon_{P_b}(A|_{W_{\mathbb R}^\perp},A|_{W_{\mathbb R}^\perp})-\{\hbar F_1^{(1)}(P_b),O\}$ with $\Upsilon_{P_b}=-\tfrac13$ units times $\{S^{(1)}(P_b),O\}$:
\begin{equation}
P_1(O)=-\big(\tfrac13+\theta_P\big)\frac{\hbar\operatorname{tr}_{\mathfrak g}(\mathrm{ad}_X\mathrm{ad}_Y)}{16\pi^2(X,Y)}\,\{S^{(1)}(P_b),O\},\qquad \{S^{(1)}(P_b),O\}=\int\omega\,\partial_b(\iota_{w_1}A,\iota_{w_2}A).
\label{eq:c9-P1}
\end{equation}
The term $-\tfrac13$ units is the limit of the cycles through $O$ that contain $S^{(1)}(P_b)$ as the composite kernel and one special edge.
\item (iii) The arity-two equations along $(\xi,P_b)$, $\xi\in\Sigma$, give $\theta_P=\theta_{P,\lambda}=\theta_{P,V}=:t_P$, and those along $(\xi_+,\xi_-)$, $\xi_\pm$ of the two chiralities, give $t_P=t_{\xi_+}+t_{\xi_-}-\tfrac{27}4-2\log\tfrac32$ (Theorem \ref{thm:c11-oddar2}). In the translation-covariant scheme of Definition \ref{def:c11-free}, $t_P=\theta=-\tfrac13$, and $P_1(O)=0$.
\end{itemize}
\end{theorem}

\begin{proof}
(i) The defect cycles along $P_b$ are those of Proposition \ref{prop:c9-structure} (ii) with $\partial_b$ in place of $v.$; every edge, vertex and counterterm commutes with $\partial_b$ (Corollary \ref{cor:c7-equiv}), so the cycles with $S^{(1)}(P_b)$ at the special edge are the translations of the arity-zero special-edge cycles, and the sum over the two-leg cycles, evaluated exactly, vanishes, in agreement with the position part of Theorem \ref{thm:c9-quadratic} (ii) (Proposition \ref{prop:c11-position} (i)). The closed and exact directions are the position parts of those of Theorem \ref{thm:c9-quadratic} (iii) (Proposition \ref{prop:c11-position}). (ii) As in Theorem \ref{thm:c9-D1}, with $\partial_b$ in place of $v.$: the graphs with $S^{(1)}(P_b)$ on an edge or on a leg reproduce the flow of the classical term $\{S^{(1)}(P_b),O\}$, the remaining graphs are the cycles through $O$ with the composite kernel and one special edge, whose limit on this part is $-\tfrac13$ units times $\{S^{(1)}(P_b),O\}$, and $-\{\hbar F_1^{(1)}(P_b),O\}$ gives $-\theta_P$ units times it. In the second evaluation of Remark \ref{rem:c9-twoeval} the composite triangle along $P_b$, symmetrised over the two legs $A|_{W_{\mathbb R}^\perp}$, gives $-2$ units and the composite two-point functional along $P_b$ on a leg of $O$ gives $\tfrac53$ units times $\{S^{(1)}(P_b),O\}$, with the same sum. (iii) Theorem \ref{thm:c11-oddar2}; in the translation-covariant scheme the position part of $\hbar F_1^{(1)}(\Delta)$ along $b$ is $2\hbar F_1^{(1)}(P_b)$, so $\theta_P=\theta$, and Theorem \ref{thm:c9-theta2} gives $\theta=-\tfrac13$.
\end{proof}
